\documentclass[11pt,a4paper]{article}
\usepackage[T1]{fontenc}
\usepackage{lmodern,amsmath,amssymb,amsthm,mathtools,microtype}
\usepackage[margin=27mm]{geometry}
\usepackage[colorlinks=true]{hyperref}
\setlength{\emergencystretch}{2em}
\newtheorem{theorem}{Theorem}[section]
\newtheorem{proposition}[theorem]{Proposition}
\newcommand{\Q}{\mathbb Q}
\title{A weight-dependent cubic trace test}
\author{}\date{17 November 2026}
\begin{document}\maketitle
\begin{abstract}
The class $1+w(2-V)$ preserves the trace pole test, its local preparation,
and its finite-algebra coordinate count. Its cubic square equation has
two rational sections and a rational point of order two. We give the
exact group-law model. Simultaneous rational realization at actual
detector traces remains unproved; the result is a new sufficient
criterion, not a smaller complete defining polynomial.
\end{abstract}

\section{Soundness, preparation and encoding}
Let $S$ contain two, let $w,V\in\Q^\times$, and assume that every
odd valuation of $w$ outside $S$ is positive. Set
\[
 b=4w/V^2,\qquad C=V^2+w(2-V),\qquad D=1+w(2-V).
\]
\begin{theorem}
If $DC$ is a rational square, then $b$ has no negative odd valuation
outside $S$. On the prepared real domain $w>0$, $0<V\le2$, both
$D$ and $C$ are strictly positive. The new class is a nonzero local
square at every fixed preparation place and every fresh selected
comparison contact of the retained general and conditional sources.
It adds no free coordinate or lost-contact height budget.
\end{theorem}
\begin{proof}
At a hypothetical negative odd pole, put $\nu=v_p(w)>0$ odd and
$m=v_p(V)>\nu/2$. Then $v_p(C)=\nu$, while
$D\equiv1\pmod p$. The product has odd valuation and is not a
square. On the real prepared domain, $D\ge1$ and $C\ge V^2>0$.

At a selected contact, $w$ has positive valuation and $V$ is integral,
so $D\in1+p\mathbb Z_p$. At a fixed preparation place, $V$ may be
taken sufficiently close to two while $w$ remains in its prescribed
bounded neighborhood; in the conditional construction $w$ tends to
zero. Thus $D$ is in the standard local square neighborhood of one,
including $1+8\mathbb Z_2$ at two. The original preparation already
makes $C$ locally square. All denominators are retained weight and
trace denominators. Strict positivity at every prepared character
makes the cleared finite square extensions totally real and etale.
Their norm guards and one-coordinate separator are the same as for
the trace-conic criterion.
\end{proof}

Thus simultaneous rational squareness of the two actual values $D_jC_j$,
with the retained point tuple, common slope and every guard, suffices for
the twenty-one- and nine-variable criteria. This simultaneous property
is an additional arithmetic requirement and is not established here.

\section{The elliptic curve and its sections}
For fixed $w\ne0$, write
\[
 \mathcal E_w:\quad
 Y^2=(1+w(2-V))(V^2+w(2-V)).
\]
Its affine cubic is
\[
 -wV^3+(w+1)^2V^2-w(4w+1)V+2w(2w+1).
\]
The change
\[
 x=1+w(2-V),\qquad z=wY
\]
gives the Weierstrass model
\begin{equation}\label{eq:model}
 z^2=x\bigl(x^2+(w^2-4w-2)x+(w+1)(3w+1)\bigr).
\end{equation}
Its discriminant is
\[
 16w^3(w-8)(w+1)^2(3w+1)^2.
\]
In particular every prepared $0<w\le1/4$ gives a nonsingular curve.
The cubic discriminant in the coordinate $V$ is
$w(w-8)(w+1)^2(3w+1)^2$.

\begin{proposition}
The points
\[
 Q=(V,Y)=(1,1+w),\qquad R=(2,2)
\]
are rational sections, and $R=-[2]Q$ with the point at infinity as
origin. The point $(x,z)=(0,0)$ in \eqref{eq:model} has order two.
In trace coordinates translation by it sends
\[
 V\longmapsto\frac{w-(1+2w)V}{1+w(2-V)}.
\]
\end{proposition}
\begin{proof}
The section equations follow by substitution. If $f(V)$ is the cubic,
then
\[
 f(V)-\bigl((1-w)V+2w\bigr)^2
       =-w(V-2)(V-1)^2.
\]
Thus the line through $Q$ and $R$ is tangent at $Q$, proving the
group-law assertion. For a model $z^2=x(x^2+ax+b)$, translation
by $(0,0)$ sends $x$ to $b/x$. Substitution of the displayed
coordinate change gives the stated formula.
\end{proof}

The section at $V=-1$, with $Y=1+3w$, is a torsion translate of
$Q$ up to sign. These sections supply rational traces on the auxiliary
curve; they do not identify those traces with the detector trace.
The rational normalization $w=u/(1+u)^2$ may be substituted in
\eqref{eq:model}, but that base change by itself gives no actual-source
incidence theorem.

\section{The linear class $V+2$ and the existing prime selection}
The class $V+2$ is positive on the prepared domain and is a unit,
congruent to two, at every hypothetical bad pole outside $S$.
At the reference trace it is four. At a selected contact with trace
residue one, it is three.

The required local square condition at that contact follows from
the already retained prime selection: the repository's
\emph{Arithmetic dependencies and whole-system completeness}, Lemma
on arbitrary avoidance, chooses $p\equiv1\pmod{8\cdot3\cdot17}$.
This is also a direct consequence of Sun's Theorem~4.1 in
arXiv:2610.20691v3: the Gaussian generators satisfy
$\pi\equiv9\pmod{16}$ and $\pi\equiv i\pmod M$, where $M$
may contain three. Taking norms gives $p\equiv1\pmod{24}$,
so quadratic reciprocity gives $(3/p)=1$. The theorem retains its
self-order and cross-order conditions and arbitrary fixed denominator
divisibility with this same $M$. Thus no extra splitting premise or
height support is required for the linear class.

\begin{thebibliography}{9}
\bibitem{Sun} Z.-W. Sun,
\emph{48 unknowns for Hilbert's tenth problem over the rationals},
arXiv:2610.20691v3 (11 November 2026), Theorem 4.1, p.~16.
\url{https://arxiv.org/pdf/2610.20691}.
\bibitem{Trace} \emph{A rational conic criterion using the actual
detector trace}, repository research manuscript, 17 November 2026.
\end{thebibliography}
\end{document}
